
\documentclass[11pt,a4paper]{ctexart}
\usepackage{geometry}
\geometry{left=2.2cm,right=2.2cm,top=2.4cm,bottom=2.4cm}
\usepackage{amsmath,amssymb,bm}
\usepackage{booktabs,longtable,array,multirow}
\usepackage{enumitem}
\usepackage{hyperref}
\usepackage{xcolor}
\usepackage{tikz}
\usetikzlibrary{arrows.meta,positioning,shapes.geometric,fit}
\usepackage{listings}
\usepackage{tcolorbox}
\tcbuselibrary{breakable,skins}
\hypersetup{colorlinks=true,linkcolor=blue!55!black,urlcolor=blue!55!black,citecolor=blue!55!black}
\setmainfont{Noto Serif}
\setCJKmainfont{Noto Sans CJK SC}
\setCJKsansfont{Noto Sans CJK SC}
\setmonofont{DejaVu Sans Mono}
\linespread{1.18}
\setlist[itemize]{leftmargin=2em,itemsep=0.25em,topsep=0.25em}
\setlist[enumerate]{leftmargin=2em,itemsep=0.25em,topsep=0.25em}
\definecolor{mainblue}{RGB}{31,70,120}
\definecolor{lightblue}{RGB}{235,242,252}
\definecolor{lightgray}{RGB}{246,246,246}
\definecolor{codebg}{RGB}{248,248,248}
\lstset{
  language=Python,
  basicstyle=\ttfamily\small,
  keywordstyle=\color{blue!70!black}\bfseries,
  commentstyle=\color{green!35!black},
  stringstyle=\color{red!55!black},
  backgroundcolor=\color{codebg},
  frame=single,
  framerule=0.3pt,
  rulecolor=\color{gray!45},
  breaklines=true,
  showstringspaces=false,
  columns=fullflexible,
  keepspaces=true,
  numbers=left,
  numberstyle=\tiny\color{gray!70}
}
\newtcolorbox{paperbox}[1]{breakable,enhanced,colback=lightblue,colframe=mainblue,title=#1,fonttitle=\bfseries,boxrule=0.8pt,arc=2mm}
\newtcolorbox{notebox}[1]{breakable,enhanced,colback=lightgray,colframe=gray!70,title=#1,fonttitle=\bfseries,boxrule=0.5pt,arc=1mm}
\newcommand{\R}{\mathbb{R}}
\newcommand{\Z}{\mathbb{Z}}
\newcommand{\argmin}{\operatorname*{arg\,min}}
\newcommand{\argmax}{\operatorname*{arg\,max}}

\title{\bfseries 规划模型建模规范\\[4pt]\large 线性规划、整数规划、0--1规划、蒙特卡洛、非线性规划、目标规划与现代优化方法}
\author{适用于数学建模论文与代码实现}
\date{\today}

\begin{document}
\maketitle

\begin{abstract}
规划模型的核心不是“评价谁更好”，而是在一组约束条件下选择决策变量，使目标函数达到最优或尽量满足多个目标。本文把常见规划模型整理为统一套路：先把实际问题转化为决策变量、目标函数和约束条件，再根据目标和约束的形式选择线性规划、整数规划、非线性规划、目标规划或随机模拟方法。文档包含基本概念、适用条件、数学原理、建模步骤、论文结论评价模板、扩展方法和 Python 代码样式。
\end{abstract}

\tableofcontents
\newpage

\section{规划模型的总主旨}

\begin{paperbox}{一句话概括}
规划问题的主旨是：在资源、规则、容量、时间、逻辑、风险等限制条件下，寻找一组可行决策变量，使某个目标最大、最小，或使多个目标尽可能达到要求。
\end{paperbox}

和 TOPSIS、FCE、熵权法、灰色关联等综合评价模型相比，规划模型更强调“\textbf{决策}”。评价模型通常回答“哪个对象更好”；规划模型通常回答“应该怎么分配、怎么安排、怎么选择”。因此它的论文写法要围绕以下三件事展开：
\[
\boxed{\text{决策变量}}\quad \boxed{\text{目标函数}}\quad \boxed{\text{约束条件}}.
\]

一般规划模型可写为
\begin{equation}
\begin{aligned}
\min_{x}\quad & f(x)\\
\text{s.t.}\quad & g_i(x)\le 0,\quad i=1,\ldots,m,\\
& h_j(x)=0,\quad j=1,\ldots,p,\\
& x\in X.
\end{aligned}
\end{equation}
其中 $x$ 是决策变量，$f(x)$ 是目标函数，$g_i,h_j$ 是约束，$X$ 表示变量类型限制，例如连续变量、整数变量、0--1变量等。

\begin{center}
\begin{tikzpicture}[node distance=1.35cm,>=Latex,font=\small]
\tikzstyle{box}=[rectangle,rounded corners,draw=mainblue,fill=lightblue,thick,align=center,minimum width=3.2cm,minimum height=0.9cm]
\node[box] (real) {现实问题\\资源分配/路径/生产/投资};
\node[box,right=of real] (data) {数据与参数\\成本、需求、容量、概率};
\node[box,right=of data] (model) {数学模型\\变量、目标、约束};
\node[box,below=of model] (solver) {求解算法\\LP/IP/NLP/仿真/启发式};
\node[box,left=of solver] (result) {结果解释\\最优值、方案、敏感性};
\node[box,left=of result] (paper) {论文输出\\模型评价、改进建议};
\draw[->,thick] (real)--(data);
\draw[->,thick] (data)--(model);
\draw[->,thick] (model)--(solver);
\draw[->,thick] (solver)--(result);
\draw[->,thick] (result)--(paper);
\end{tikzpicture}
\end{center}

\subsection{论文中的总建模步骤}
\begin{enumerate}
\item \textbf{明确决策目标。} 是最小成本、最大收益、最短时间、最低风险，还是多个目标共同满足？
\item \textbf{定义决策变量。} 变量是连续数量、整数数量、是否选择、路径顺序，还是随机样本？
\item \textbf{建立目标函数。} 把“最好”写成数学表达式，例如 $\min c^Tx$ 或 $\max p^Tx$。
\item \textbf{写出约束条件。} 资源、容量、时间、逻辑、供需平衡、预算、政策要求都要转成等式或不等式。
\item \textbf{判断模型类别。} 线性、整数、0--1、非线性、随机、目标规划、多目标规划等。
\item \textbf{选择求解方法。} 线性规划可用单纯形法或内点法，整数规划可用分支定界/割平面/CP-SAT，非线性规划可用梯度、牛顿、SLSQP、trust-constr 或全局优化方法。
\item \textbf{解释结果。} 包括最优方案、最优值、约束是否紧、敏感性分析、模型局限和改进建议。
\end{enumerate}

\section{规划模型分类总表}

\begin{longtable}{p{2.4cm}p{3.2cm}p{4.8cm}p{4.3cm}}
\toprule
模型类别 & 标准形式 & 适合问题 & 常见方法 \\
\midrule
线性规划 LP & $\min c^Tx,\ Ax\le b$ & 资源分配、生产计划、运输、配料、投资组合线性近似 & 单纯形法、内点法、SciPy linprog、OR-Tools、Pyomo \\
整数规划 IP & LP 加 $x_j\in\Z$ & 人数、台数、车辆数、批次数等必须为整数 & 分支定界、割平面、分支割、CP-SAT \\
0--1规划 BIP & $x_j\in\{0,1\}$ & 选或不选、开或不开、是否分配、路径边是否采用 & 大 $M$ 法、逻辑约束、背包、指派、选址、CP-SAT \\
蒙特卡洛法 & 随机抽样估计 $E[g(X)]$ & 不确定性模拟、风险分析、随机参数、难以解析的问题 & 随机采样、情景模拟、置信区间、随机搜索 \\
非线性规划 NLP & $\min f(x)$, $f,g,h$ 可非线性 & 成本曲线、收益递减、物理模型、机器学习参数 & 梯度法、牛顿法、拉格朗日、KKT、SLSQP、trust-constr \\
目标规划 GP & 最小化目标偏差变量 & 多目标同时满足、目标有期望值、指标有优先级 & 加权目标规划、优先等级法、Chebyshev 目标规划 \\
凸优化 & $f,g_i$ 凸，$h_j$ 仿射 & 大规模可靠求解、稀疏回归、投资组合、控制 & CVXPY、内点法、ADMM、近端梯度 \\
随机/鲁棒优化 & 参数不确定 & 需求、价格、天气、风险不确定 & 情景规划、机会约束、鲁棒集、两阶段规划 \\
启发式/元启发式 & 不要求解析可导 & 旅行商、排班、路径、复杂组合优化 & 遗传算法、模拟退火、粒子群、差分进化、贝叶斯优化 \\
\bottomrule
\end{longtable}

\section{线性规划}

\subsection{概念与标准形式}
线性规划是最基础的规划模型。若目标函数和约束条件都是决策变量的一次函数，则可写成：
\begin{equation}
\begin{aligned}
\min\quad & c^Tx\\
\text{s.t.}\quad & Ax\le b,\\
& A_{eq}x=b_{eq},\\
& l\le x\le u.
\end{aligned}
\end{equation}
若题目是最大化，可以通过 $\max z= c^Tx$ 转化为 $\min -c^Tx$。

\subsection{使用条件}
\begin{itemize}
\item \textbf{线性比例性：} 每增加一单位变量，对目标和约束的影响是固定的。
\item \textbf{可加性：} 总成本、总收益、总资源消耗可以分项相加。
\item \textbf{确定性：} 参数 $c,A,b$ 在模型期内视为已知。
\item \textbf{连续性：} 变量可取连续值。若必须为整数，就转为整数规划。
\end{itemize}

\subsection{建模模板}
\begin{enumerate}
\item 设变量 $x_j$ 表示第 $j$ 类活动的数量。
\item 写目标：最小化总成本 $\min \sum_j c_jx_j$ 或最大化总收益 $\max \sum_j p_jx_j$。
\item 写约束：资源限制 $\sum_j a_{ij}x_j\le b_i$，需求约束 $\sum_j d_{ij}x_j\ge q_i$。
\item 加非负约束 $x_j\ge0$。
\item 求解后说明哪些约束是紧约束，哪些资源有剩余。
\end{enumerate}

\subsection{Python 代码样式：SciPy linprog}
\begin{lstlisting}
import numpy as np
from scipy.optimize import linprog

# min 3x1 + 5x2
# s.t. 2x1 + x2 <= 100
#      x1 + 2x2 <= 80
#      x1, x2 >= 0
c = np.array([3, 5])
A_ub = np.array([[2, 1], [1, 2]])
b_ub = np.array([100, 80])
bounds = [(0, None), (0, None)]

res = linprog(c, A_ub=A_ub, b_ub=b_ub, bounds=bounds, method="highs")
print(res.success, res.fun, res.x)
\end{lstlisting}

\section{整数规划与 0--1 规划}

\subsection{整数规划}
整数规划在线性规划基础上增加整数约束：
\begin{equation}
\begin{aligned}
\min\quad & c^Tx\\
\text{s.t.}\quad & Ax\le b,\\
& x_j\in \Z,\quad j\in I.
\end{aligned}
\end{equation}
如果只有部分变量要求整数，称为混合整数规划（MIP/MILP）。

整数规划适合变量本身不能拆分的情形，例如机器台数、员工人数、车辆数、课程安排、货物批次数。它通常比线性规划难得多，因为可行域不再是连续多面体，而是许多离散点。

\subsection{0--1 规划}
0--1 规划是整数规划的特殊形式：
\[
 x_j\in\{0,1\}.
\]
它最适合表达“是否”类决策：是否选择项目、是否建设仓库、是否把任务分配给某人、路径中是否经过某条边。

\subsection{常用建模技巧}
\begin{paperbox}{0--1 变量的常见含义}
\begin{itemize}
\item $x_j=1$ 表示选择第 $j$ 个项目，$x_j=0$ 表示不选择。
\item $y_{ij}=1$ 表示任务 $i$ 分配给人员 $j$，否则为 $0$。
\item $z_k=1$ 表示第 $k$ 个仓库开启，否则关闭。
\end{itemize}
\end{paperbox}

\begin{enumerate}
\item \textbf{背包模型：}
\[
\max \sum_j v_jx_j,\quad \sum_j w_jx_j\le W,\quad x_j\in\{0,1\}.
\]
\item \textbf{指派模型：}
\[
\min\sum_i\sum_j c_{ij}x_{ij},\quad \sum_j x_{ij}=1,\quad \sum_i x_{ij}=1.
\]
\item \textbf{大 $M$ 法：} 若 $y=0$ 时要求 $x=0$，$y=1$ 时允许 $x$ 取正，可写
\[
0\le x\le My,\quad y\in\{0,1\}.
\]
\item \textbf{逻辑约束：} 若“选择 A 则必须选择 B”，可写 $x_A\le x_B$。
\end{enumerate}

\subsection{Python 代码样式：OR-Tools CP-SAT 0--1 背包}
\begin{lstlisting}
from ortools.sat.python import cp_model

values = [10, 8, 15, 4]
weights = [5, 4, 8, 3]
capacity = 12
n = len(values)

model = cp_model.CpModel()
x = [model.NewBoolVar(f"x{j}") for j in range(n)]

model.Add(sum(weights[j] * x[j] for j in range(n)) <= capacity)
model.Maximize(sum(values[j] * x[j] for j in range(n)))

solver = cp_model.CpSolver()
status = solver.Solve(model)

if status in (cp_model.OPTIMAL, cp_model.FEASIBLE):
    print("objective =", solver.ObjectiveValue())
    print([solver.Value(x[j]) for j in range(n)])
\end{lstlisting}

\section{蒙特卡洛法}

\subsection{它是不是规划模型？}
蒙特卡洛法本身不是传统意义上的“规划模型”，而是一种随机模拟方法。它经常与规划模型结合，用来处理参数不确定、风险评估或复杂模型无法解析求解的情况。其主旨是：用大量随机样本近似期望、概率、风险或目标函数分布。

例如需求 $D$、价格 $P$、成本 $C$ 不确定时，可以设定它们的概率分布，重复抽样并计算利润：
\[
\Pi = P\cdot \min(Q,D)-C Q.
\]
通过模拟得到 $E(\Pi)$、亏损概率、最坏 5\% 情况等。

\subsection{建模步骤}
\begin{enumerate}
\item 确定随机变量，例如需求、价格、故障率、服务时间。
\item 选择概率分布，例如正态、三角、均匀、泊松、经验分布。
\item 生成大量随机样本。
\item 对每个样本计算目标值或约束是否满足。
\item 得到均值、方差、置信区间、风险概率和敏感性排序。
\item 如果需要优化，则在不同决策方案下重复模拟，选择期望收益高或风险小的方案。
\end{enumerate}

\subsection{Python 代码样式：库存订货的模拟优化}
\begin{lstlisting}
import numpy as np

rng = np.random.default_rng(0)
N = 20000
price = 12
cost = 7
salvage = 2

# demand follows a truncated normal-like sample
demand = np.maximum(0, rng.normal(loc=100, scale=20, size=N))

def simulate_profit(Q):
    sales = np.minimum(Q, demand)
    leftover = np.maximum(Q - demand, 0)
    profit = price * sales + salvage * leftover - cost * Q
    return profit.mean(), np.percentile(profit, 5)

candidates = np.arange(50, 151)
records = [(Q, *simulate_profit(Q)) for Q in candidates]
best = max(records, key=lambda t: t[1])
print("best by expected profit:", best)
\end{lstlisting}

\section{非线性规划}

\subsection{无约束极值}
无约束非线性规划写为
\[
\min_{x\in\R^n} f(x).
\]
若 $f$ 可导，则一阶必要条件为
\[
\nabla f(x^*)=0.
\]
若 Hessian 矩阵 $\nabla^2 f(x^*)$ 正定，则 $x^*$ 是严格局部极小点；若负定，则是局部极大点；若不定，则可能是鞍点。

\subsection{等式约束：拉格朗日乘子法}
若问题为
\[
\min f(x),\quad h_i(x)=0,
\]
构造拉格朗日函数
\[
L(x,\lambda)=f(x)+\sum_i\lambda_i h_i(x).
\]
候选极值点满足
\[
\nabla_x L(x,\lambda)=0,\quad h_i(x)=0.
\]
拉格朗日乘子 $\lambda_i$ 可以理解为约束右端项变化时目标最优值的边际变化。

\subsection{不等式约束：KKT 条件}
若问题为
\[
\min f(x),\quad g_i(x)\le0,\quad h_j(x)=0,
\]
KKT 条件为
\begin{align}
&\nabla f(x^*)+\sum_i\mu_i\nabla g_i(x^*)+\sum_j\lambda_j\nabla h_j(x^*)=0,\\
&g_i(x^*)\le0,
\quad h_j(x^*)=0,\\
&\mu_i\ge0,\\
&\mu_i g_i(x^*)=0.
\end{align}
其中 $\mu_i g_i(x^*)=0$ 叫互补松弛：一个不等式约束要么紧，要么对应乘子为 0。

\subsection{Python 代码样式：SLSQP 求有约束非线性规划}
\begin{lstlisting}
import numpy as np
from scipy.optimize import minimize

# min (x-1)^2 + (y-2)^2
# s.t. x + y >= 2, x^2 + y^2 <= 10

def obj(v):
    x, y = v
    return (x - 1)**2 + (y - 2)**2

constraints = [
    {"type": "ineq", "fun": lambda v: v[0] + v[1] - 2},
    {"type": "ineq", "fun": lambda v: 10 - v[0]**2 - v[1]**2},
]
res = minimize(obj, x0=np.array([0.0, 0.0]), constraints=constraints,
               method="SLSQP")
print(res.success, res.fun, res.x)
\end{lstlisting}

\section{目标规划}

\subsection{为什么需要目标规划}
实际问题常常不只有一个目标。例如企业希望成本低、产量高、污染少、员工加班少。若强行把所有目标合成一个目标，权重很难解释；若逐个优化，目标之间又可能冲突。目标规划的思想是：给每个目标设置期望值，再最小化对目标的偏离程度。

设第 $k$ 个目标的期望值为 $G_k$，目标表达式为 $f_k(x)$，引入正偏差 $d_k^+$ 和负偏差 $d_k^-$：
\begin{equation}
f_k(x)+d_k^- - d_k^+ = G_k,
\quad d_k^+,d_k^-\ge0.
\end{equation}
其中 $d_k^-$ 表示未达到目标的不足量，$d_k^+$ 表示超过目标的超出量。根据实际含义，只惩罚不希望出现的偏差。

\subsection{加权系数法}
把不同目标偏差加权求和：
\begin{equation}
\min Z=\sum_k \alpha_k d_k^- + \beta_k d_k^+.
\end{equation}
适合目标优先级差别不绝对、可以进行权衡的情况。

\subsection{优先等级法}
若目标有绝对优先级，例如安全目标高于成本目标，高优先级目标的任何改善都比低优先级目标重要，则使用优先等级法。可按优先级分层：
\[
P_1\gg P_2\gg P_3.
\]
求解时通常先最小化第一级偏差，固定第一级最优值后，再优化第二级偏差，以此类推。这也叫词典序或预emptive 目标规划。

\subsection{Chebyshev 目标规划}
若希望避免某个目标偏差特别大，可最小化最大归一化偏差：
\begin{equation}
\min t,\quad \frac{d_k}{s_k}\le t,\quad k=1,\ldots,K.
\end{equation}
它强调目标之间的均衡。

\subsection{Python 代码样式：加权目标规划}
\begin{lstlisting}
import numpy as np
from scipy.optimize import linprog

# Example: production x1, x2.
# Goals: profit >= 100, labor <= 40.
# profit = 6x1 + 8x2, labor = 2x1 + 3x2.
# Variables: x1,x2,d1_minus,d1_plus,d2_minus,d2_plus

c = np.array([0, 0, 5, 0, 0, 2])  # penalize profit shortfall and labor excess

A_eq = np.array([
    [6, 8, 1, -1, 0, 0],   # profit + d- - d+ = 100
    [2, 3, 0, 0, 1, -1],   # labor + d- - d+ = 40
])
b_eq = np.array([100, 40])

bounds = [(0, None)] * 6
res = linprog(c, A_eq=A_eq, b_eq=b_eq, bounds=bounds, method="highs")
print(res.success, res.fun, res.x)
\end{lstlisting}

\section{现代常用方法与补充}

\subsection{凸优化}
如果目标函数是凸函数，约束集合是凸集，则局部最优就是全局最优。许多看似非线性的模型其实属于凸优化，例如最小二乘、岭回归、Lasso、二次规划、二阶锥规划、部分投资组合模型。凸优化的优点是理论和算法都比较成熟，求解可靠性强。

\subsection{二次规划 QP}
二次规划常见形式为
\[
\min \frac12x^TQx+c^Tx,
\quad Ax\le b.
\]
若 $Q\succeq0$，则为凸二次规划。它常用于投资组合、最小方差、轨迹平滑、支持向量机等。

\subsection{动态规划}
动态规划适合多阶段决策问题。其核心是状态转移和 Bellman 最优性原理：
\[
V_t(s)=\min_a\{c_t(s,a)+V_{t+1}(s')\}.
\]
适用于库存控制、最短路径、设备更新、背包、资源分配等问题。

\subsection{网络流与路径优化}
最短路、最大流、最小费用流、运输问题、车辆路径问题都是规划模型的重要分支。若题目中出现“节点、边、路径、流量、容量”，应优先考虑网络优化模型。

\subsection{随机规划与鲁棒优化}
随机规划把不确定参数表示为随机变量或情景集合，通常优化期望成本或风险指标。鲁棒优化不假设准确概率分布，而是假设参数落在不确定集合中，并要求方案在最坏情况下仍可行或表现较好。

\subsection{启发式与元启发式}
当模型过于复杂、变量很多、非凸、离散、难以获得精确最优解时，可以使用启发式方法：
\begin{itemize}
\item 遗传算法 GA：适合编码清楚、可行解容易生成的问题。
\item 模拟退火 SA：适合跳出局部最优。
\item 粒子群 PSO：适合连续变量黑箱优化。
\item 差分进化 DE：适合边界已知的全局优化。
\item 贝叶斯优化：适合每次评估代价高的黑箱函数，例如机器学习超参数。
\end{itemize}

\subsection{机器学习与优化结合}
现代应用中经常出现“预测 + 优化”。先用机器学习预测需求、价格、风险、时间，再把预测值放入规划模型。更进一步的 end-to-end predict-then-optimize 会直接以最终决策质量训练预测模型。建模论文中若时间有限，可写成扩展思路，不必全部实现。

\section{如何根据题目选择模型}

\begin{longtable}{p{4cm}p{4cm}p{6.2cm}}
\toprule
题目关键词 & 优先模型 & 判断依据 \\
\midrule
成本最低、收益最大、资源限制、变量连续 & 线性规划 & 目标和约束都是一次式 \\
人数、台数、批次、班次必须整数 & 整数规划 & 变量不可拆分 \\
选或不选、开或不开、是否分配 & 0--1规划 & 逻辑选择和组合决策 \\
路径、车辆、配送、旅行商 & 网络优化/车辆路径/整数规划 & 节点和边结构明显 \\
风险、不确定、随机需求 & 蒙特卡洛/随机规划/鲁棒优化 & 参数不是固定数 \\
收益递减、成本曲线、非线性物理关系 & 非线性规划 & 目标或约束含乘积、幂、指数、对数等 \\
多个目标有期望值 & 目标规划 & 目标不一定最优，而是尽量接近目标值 \\
目标之间有绝对优先级 & 优先等级目标规划 & 高优先级目标不可牺牲 \\
函数黑箱、不可导、局部最优多 & 元启发式/全局优化 & 解析条件难写，数值评估可行 \\
模型规模大但凸 & 凸优化/CVXPY & 可用标准凸优化工具可靠求解 \\
\bottomrule
\end{longtable}

\section{论文写作模板}

\subsection{模型建立段落}
\begin{paperbox}{可直接改写}
为解决该资源配置问题，本文建立规划模型。设 $x_j$ 表示第 $j$ 类决策活动的实施水平，根据问题目标构造目标函数 $f(x)$，并将预算、容量、需求、逻辑关系和变量取值范围转化为约束条件。若目标函数和约束均为线性形式，则采用线性规划；若部分变量具有整数或是否选择含义，则进一步建立整数规划或 0--1 规划。模型求解后，根据最优解、目标函数值、约束紧性和敏感性分析评价方案的可行性与稳定性。
\end{paperbox}

\subsection{结果分析段落}
\begin{paperbox}{结论评价模板}
由求解结果可知，模型在满足全部约束条件的前提下取得最优目标值 $Z^*$，对应决策方案为 $x^*$。从变量取值看，资源主要分配给了……，说明这些活动在目标函数中具有较高收益或较低成本。从约束松弛量看，……约束达到边界，属于限制系统性能的关键约束；……资源仍有剩余，对最优值影响较小。若对关键参数进行扰动，最优方案变化不大，说明模型结论具有一定稳定性。
\end{paperbox}

\subsection{局限性段落}
\begin{notebox}{常见局限写法}
本文模型将部分参数视为确定值，未完全考虑需求波动、价格变化和执行过程中的随机扰动；同时，线性化处理可能简化了真实的非线性关系。后续可引入随机规划、鲁棒优化或蒙特卡洛模拟，对不确定性下的方案表现进行进一步检验。
\end{notebox}

\section{完整代码模板}
以下代码把常用规划模型的写法放在一个文件中。实际使用时不需要全用，只选与题目匹配的部分。

\begin{lstlisting}
import numpy as np
from scipy.optimize import linprog, minimize, differential_evolution

# ---------- Linear Programming ----------
def solve_lp():
    c = np.array([3, 5])
    A_ub = np.array([[2, 1], [1, 2]])
    b_ub = np.array([100, 80])
    res = linprog(c, A_ub=A_ub, b_ub=b_ub,
                  bounds=[(0, None), (0, None)], method="highs")
    return res

# ---------- Nonlinear Programming ----------
def solve_nlp():
    def f(v):
        x, y = v
        return (x - 1)**2 + (y - 2)**2
    cons = [
        {"type": "ineq", "fun": lambda v: v[0] + v[1] - 2},
        {"type": "ineq", "fun": lambda v: 10 - v[0]**2 - v[1]**2},
    ]
    return minimize(f, x0=np.array([0.0, 0.0]), constraints=cons, method="SLSQP")

# ---------- Monte Carlo Simulation ----------
def monte_carlo_inventory(seed=0):
    rng = np.random.default_rng(seed)
    demand = np.maximum(0, rng.normal(100, 20, size=20000))
    price, cost, salvage = 12, 7, 2
    def profit(Q):
        sales = np.minimum(Q, demand)
        leftover = np.maximum(Q - demand, 0)
        p = price * sales + salvage * leftover - cost * Q
        return p.mean(), np.percentile(p, 5)
    qs = np.arange(50, 151)
    records = [(Q, *profit(Q)) for Q in qs]
    return max(records, key=lambda r: r[1])

# ---------- Global Optimization ----------
def global_optimize():
    def rastrigin(x):
        x = np.asarray(x)
        return 20 + np.sum(x**2 - 10 * np.cos(2 * np.pi * x))
    bounds = [(-5.12, 5.12), (-5.12, 5.12)]
    return differential_evolution(rastrigin, bounds, seed=0)

if __name__ == "__main__":
    print("LP:", solve_lp().fun, solve_lp().x)
    print("NLP:", solve_nlp().fun, solve_nlp().x)
    print("Monte Carlo:", monte_carlo_inventory())
    print("Global:", global_optimize().fun, global_optimize().x)
\end{lstlisting}

\section{OR-Tools 0--1 规划代码}
\begin{lstlisting}
from ortools.sat.python import cp_model

def solve_binary_knapsack(values, weights, capacity):
    model = cp_model.CpModel()
    n = len(values)
    x = [model.NewBoolVar(f"x{j}") for j in range(n)]
    model.Add(sum(weights[j] * x[j] for j in range(n)) <= capacity)
    model.Maximize(sum(values[j] * x[j] for j in range(n)))
    solver = cp_model.CpSolver()
    status = solver.Solve(model)
    if status not in (cp_model.OPTIMAL, cp_model.FEASIBLE):
        return None
    return {
        "objective": solver.ObjectiveValue(),
        "x": [solver.Value(var) for var in x],
        "status": solver.StatusName(status),
    }
\end{lstlisting}

\section{CVXPY 凸优化代码}
\begin{lstlisting}
import cvxpy as cp
import numpy as np

# Portfolio-like convex quadratic programming
mu = np.array([0.08, 0.12, 0.10])
Sigma = np.array([
    [0.10, 0.02, 0.01],
    [0.02, 0.15, 0.03],
    [0.01, 0.03, 0.12],
])

x = cp.Variable(3, nonneg=True)
risk_aversion = 2.0
objective = cp.Maximize(mu @ x - risk_aversion * cp.quad_form(x, Sigma))
constraints = [cp.sum(x) == 1]
prob = cp.Problem(objective, constraints)
prob.solve()
print(prob.value, x.value)
\end{lstlisting}

\section{参考资料与网页实例}
\begin{enumerate}
\item Boyd, S. and Vandenberghe, L. \textit{Convex Optimization}. Stanford University. 该书系统介绍凸优化，并说明线性规划和最小二乘是凸优化中的重要特例。
\item SciPy Documentation: \texttt{scipy.optimize}. 说明 SciPy 提供线性规划、非线性优化、全局优化、最小二乘等函数。
\item Google OR-Tools Documentation: Linear Programming and CP-SAT Solver. 官方文档给出 LP 求解基本步骤，并说明 CP-SAT 适合整数规划和约束规划。
\item NEOS Guide: Integer Programming. 给出整数规划标准形式与基本概念。
\item Pyomo Documentation. Pyomo 是 Python 优化建模语言，可表达线性、整数、非线性、混合整数非线性等问题。
\item CVXPY Documentation and Diamond \& Boyd paper. CVXPY 是嵌入 Python 的凸优化建模语言，能按数学形式表达问题并调用求解器。
\item SciPy Documentation: \texttt{differential\_evolution} and \texttt{dual\_annealing}. 用于全局优化和随机启发式求解。
\item Goal Programming resources. 目标规划是多目标优化分支，通过偏差变量、加权偏差或优先等级处理多目标冲突。
\item Optuna Documentation. Optuna 是机器学习中常用的超参数优化框架，体现了现代黑箱优化和贝叶斯优化思想。
\end{enumerate}

\end{document}
